\documentclass[aspectratio=169,10pt]{ctexbeamer}

\usepackage{graphicx}
\usepackage{booktabs}
\usepackage{tabularx}
\usepackage{colortbl}
\usepackage{ragged2e}
\usepackage{tikz}
\usepackage{tcolorbox}
\usepackage{listings}
\usepackage{xcolor}

\usetikzlibrary{arrows.meta,positioning,shapes.geometric,calc}

% ---------- color palette ----------
\definecolor{deepgreen}{HTML}{065F46}
\definecolor{green}{HTML}{059669}
\definecolor{greenlight}{HTML}{ECFDF5}
\definecolor{blue}{HTML}{0284C7}
\definecolor{bluelight}{HTML}{E0F2FE}
\definecolor{orange}{HTML}{D97706}
\definecolor{orangelight}{HTML}{FEF3C7}
\definecolor{red}{HTML}{E11D48}
\definecolor{navy}{HTML}{0F172A}
\definecolor{text}{HTML}{1E293B}
\definecolor{muted}{HTML}{64748B}
\definecolor{line}{HTML}{E2E8F0}

\setbeamercolor{normal text}{fg=text}
\setbeamercolor{frametitle}{fg=navy}
\setbeamercolor{footline}{fg=muted}
\setbeamertemplate{navigation symbols}{}
\setbeamertemplate{footline}{%
  \hspace{0.55cm}\color{muted}\scriptsize DLCPD-25 · Plan-A 模型架构\quad
  \hfill\raisebox{0.15cm}{\color{green}\rule{0.6cm}{1.2pt}}\hspace{0.2cm}
  \color{muted}\scriptsize\insertframenumber/\inserttotalframenumber\hspace{0.55cm}\vspace{0.12cm}}

\setbeamertemplate{frametitle}{%
  \vspace{0.5cm}\hspace{0.55cm}{\large\bfseries\insertframetitle}\par
  \vspace{0.12cm}\hspace{0.55cm}\color{green}\rule{0.08\paperwidth}{1.6pt}\par
  \vspace{0.28cm}}

\setlength{\parskip}{4pt}

% ---------- reusable cards ----------
\newtcolorbox{card}[2][]{%
  colback=#2, colframe=line, boxrule=0.5pt, arc=2.5mm,
  left=2mm, right=2mm, top=1.5mm, bottom=1.5mm,
  fontupper=\small, #1}

\newcommand{\carditem}[3]{%
  \begin{card}{#3}
    \textbf{#1}\\[1pt] {\color{muted}\small #2}
  \end{card}}

\newcommand{\fig}[3]{%
  \begin{center}\includegraphics[width=#1\textwidth,height=#2\textheight,keepaspectratio]{#3}\end{center}}

% ---------- document ----------
\begin{document}

% Title
{
\setbeamertemplate{footline}{}
\begin{frame}[plain]
\begin{tikzpicture}[remember picture,overlay]
  \fill[deepgreen] (current page.south west) rectangle (current page.north east);
  \fill[green] ([yshift=-0.08\paperheight]current page.north west) rectangle ([yshift=-0.085\paperheight]current page.north east);
  \fill[green] (current page.north west) rectangle ++(0.13\paperwidth,-0.015\paperheight);
  \node[text=white,anchor=west,font=\LARGE\bfseries] at ([xshift=0.09\paperwidth,yshift=-0.24\paperheight]current page.north west)
    {基于 DLCPD-25 的农产品病虫害与缺陷分类目标检测系统};
  \node[text=white,anchor=west,font=\LARGE\bfseries] at ([xshift=0.09\paperwidth,yshift=-0.32\paperheight]current page.north west)
    {模型架构详解};
  \node[text=green!30!white,anchor=west,font=\normalsize] at ([xshift=0.09\paperwidth,yshift=-0.43\paperheight]current page.north west)
    {Plan-A 双专家架构：ConvNeXt-Tiny 分类专家 + ConvNeXt-Tiny-FPN Faster R-CNN 检测专家};
  \node[text=green!30!white,anchor=west,font=\small] at ([xshift=0.09\paperwidth,yshift=-0.80\paperheight]current page.north west)
    {2026-08-18};
\end{tikzpicture}
\end{frame}
}

% Agenda
\begin{frame}{目录}
\begin{columns}[T]
\begin{column}{0.48\textwidth}
\carditem{01\quad 任务定义与数据概览}{DLCPD-25 分类 + IP102 检测}{greenlight}
\vspace{2mm}
\carditem{02\quad Plan-A 总体架构}{双专家独立训练、统一编排}{greenlight}
\vspace{2mm}
\carditem{03\quad 分类专家：ConvNeXt-Tiny}{主干结构、Block 与多任务头}{greenlight}
\vspace{2mm}
\carditem{04\quad 分类训练策略}{Focal Loss、类别平衡、EMA、早停}{greenlight}
\vspace{2mm}
\carditem{05\quad 分类训练过程与结果}{验证曲线与最终测试指标}{greenlight}
\vspace{2mm}
\carditem{06\quad 检测专家：FPN + Faster R-CNN}{RPN、RoIAlign、RoI Heads}{bluelight}
\end{column}
\begin{column}{0.48\textwidth}
\carditem{07\quad 检测训练策略}{Repeat Factor 采样、两阶段损失}{bluelight}
\vspace{2mm}
\carditem{08\quad 检测训练过程与结果}{验证 mAP 曲线与最终测试指标}{bluelight}
\vspace{2mm}
\carditem{09\quad 关键设计决策}{分辨率、骨干复用、EMA、测试隔离}{orangelight}
\vspace{2mm}
\carditem{10\quad 推理系统与 Web 应用}{分类 + 检测 + 标注图}{orangelight}
\vspace{2mm}
\carditem{11\quad 工程实现与复现}{代码结构、训练与评估命令}{greenlight}
\vspace{2mm}
\carditem{12\quad 限制与未来工作}{小目标、稀有害虫、部署优化}{greenlight}
\end{column}
\end{columns}
\end{frame}

% Tasks
\begin{frame}{任务定义}
\begin{columns}[T]
\begin{column}{0.48\textwidth}
\carditem{DLCPD-25 细粒度分类}{输入整张农作物图片，识别 203 类病虫害、健康状态与生理缺陷，输出 Top-5。标签按 22 个宿主作物与 4 类属性组织。}{greenlight}
\end{column}
\begin{column}{0.48\textwidth}
\carditem{IP102 害虫目标检测}{定位并识别 96 类农业害虫，输出边界框与置信度；检测类别映射回 DLCPD-25 公共 class\_id。}{bluelight}
\end{column}
\end{columns}
\vspace{3mm}
\carditem{系统输出合同}{Web 上传图片后同时返回：整图 Top-5 分类、96 类害虫检测框和标注图。病害、健康、生理缺陷只有图像级标签，因此不生成虚假检测框。}{orangelight}
\end{frame}

% Data
\begin{frame}{数据概览}
\begin{table}
\centering\small
\begin{tabular}{lllll}
\toprule
\textbf{数据集} & \textbf{任务} & \textbf{样本} & \textbf{类别/标注} & \textbf{关键特点} \\
\midrule
DLCPD-25 & 整图分类 & 221{,}396 张 & 203 类，22 宿主 & 单类 5$\sim$10{,}619 张，长尾严重 \\
IP102 Detection & 目标检测 & 18{,}976 张 & 96 类，22{,}283 框 & 约 93\% 单目标，目标中位面积 0.35 \\
\bottomrule
\end{tabular}
\end{table}
\vspace{1mm}
\begin{columns}[T]
\begin{column}{0.48\textwidth}\carditem{分类核心矛盾}{长尾分布导致的类别不均衡}{greenlight}\end{column}
\begin{column}{0.48\textwidth}\carditem{检测核心矛盾}{细粒度害虫外观相似与定位精度}{bluelight}\end{column}
\end{columns}
\end{frame}

% Overview image
\begin{frame}{Plan-A 总体架构}
\fig{0.95}{0.62}{assets/overview.png}
\end{frame}

% Motivation
\begin{frame}{为什么采用双专家架构}
\begin{columns}[T]
\begin{column}{0.48\textwidth}
\carditem{任务目标不同}{分类是 203 类整图识别；检测是 96 类定位 + 分类，标签空间与损失函数不同。}{greenlight}
\vspace{2mm}
\carditem{分辨率需求不同}{分类 384 保留纹理，检测 640 保留定位细节。}{bluelight}
\vspace{2mm}
\carditem{长尾处理不同}{分类类别平衡采样，检测 Repeat Factor Sampling。}{orangelight}
\end{column}
\begin{column}{0.48\textwidth}
\carditem{训练与选型独立}{学习率、早停、EMA、选择指标可独立设置。}{greenlight}
\vspace{2mm}
\carditem{部署与维护灵活}{两个专家可分别导出、替换和升级。}{bluelight}
\end{column}
\end{columns}
\end{frame}

% Classification overview
\begin{frame}{分类专家：ConvNeXt-Tiny}
\begin{columns}[T]
\begin{column}{0.52\textwidth}
\small
\begin{tabular}{p{0.28\textwidth}p{0.68\textwidth}}
\toprule
\textbf{模块} & \textbf{结构} \\
\midrule
Stem & 4$\times$4 Conv, stride 4，输出 96 通道 \\
Stage 1 & 3 个 ConvNeXt Block，96 通道 \\
Stage 2 & 下采样 + 3 Block，192 通道 \\
Stage 3 & 下采样 + 9 Block，384 通道 \\
Stage 4 & 下采样 + 3 Block，768 通道 \\
Pooling & GAP + LayerNorm，输出 768 维 \\
主分类头 & Linear(768, 203) \\
宿主辅助头 & Linear(768, 22)，仅训练时使用 \\
属性辅助头 & Linear(768, 4) \\
\bottomrule
\end{tabular}
\end{column}
\begin{column}{0.44\textwidth}
\fig{1.0}{0.46}{assets/convnext_stages.png}
\end{column}
\end{columns}
\end{frame}

% ConvNeXt block
\begin{frame}{ConvNeXt Block：核心计算单元}
\fig{0.95}{0.58}{assets/convnext_block.png}
\vspace{-2mm}
\begin{card}{greenlight}
\textbf{解读：}先用 7$\times$7 Depthwise Conv 提取空间特征，再用两个 1$\times$1 Linear 完成通道升维与降维；LayerNorm 和 Layer Scale 借鉴 Transformer 训练稳定性，残差连接保留梯度通路。
\end{card}
\end{frame}

% Classification training
\begin{frame}{分类训练策略与损失}
\begin{card}{greenlight}
\textbf{Focal Loss：} $\mathcal{L} = -\alpha_t\,(1-p_t)^{\gamma}\,\log p_t$，其中 $\gamma=2$，聚焦难样本；配合类别平衡权重抑制头部类主导。
\end{card}
\vspace{2mm}
\begin{columns}[T]
\begin{column}{0.48\textwidth}
\carditem{类别平衡采样}{按 $\sqrt{1/N_c}$ 设计采样权重，提升尾部类出现频率。}{greenlight}
\vspace{2mm}
\carditem{层级辅助监督}{22 宿主头 + 4 属性头，帮助主干学习层级结构。}{bluelight}
\end{column}
\begin{column}{0.48\textwidth}
\carditem{优化与正则}{AdamW；warmup 3 epochs + cosine；EMA 0.999；早停 patience 6。}{orangelight}
\end{column}
\end{columns}
\end{frame}

% Classification history
\begin{frame}{分类模型训练过程}
\fig{0.95}{0.62}{assets/classification_history.png}
\vspace{-2mm}
\begin{card}{greenlight}
第 5 轮验证 Macro-F1 最高（74.38\%），连续 6 轮未提升后早停；最终测试集仅评估一次。
\end{card}
\end{frame}

% Classification final metrics
\begin{frame}{分类模型最终测试结果}
\begin{columns}[T]
\begin{column}{0.5\textwidth}
\begin{table}\centering\small
\begin{tabular}{lll}
\toprule
\textbf{指标} & \textbf{结果} & \textbf{含义} \\
\midrule
Top-1 & 89.2240\% & 第一名预测正确 \\
Top-5 & 97.7366\% & 真实类别进入前五名 \\
Macro-F1 & 75.3537\% & 203 类 F1 平均 \\
Balanced Acc & 77.5782\% & 类别召回率平均 \\
测试样本 & 22{,}179 张 & 仅评估一次 \\
\bottomrule
\end{tabular}
\end{table}
\end{column}
\begin{column}{0.46\textwidth}
\fig{1.0}{0.46}{assets/classification_metrics.png}
\end{column}
\end{columns}
\end{frame}

% Detection architecture
\begin{frame}{检测专家：ConvNeXt-Tiny-FPN + Faster R-CNN}
\fig{0.95}{0.66}{assets/faster_rcnn.png}
\end{frame}

% Detection training
\begin{frame}{检测训练策略与损失}
\begin{columns}[T]
\begin{column}{0.48\textwidth}
\carditem{骨干初始化}{ConvNeXt 骨干加载分类专家 EMA 权重，检测头随机初始化。}{greenlight}
\vspace{2mm}
\carditem{Repeat Factor Sampling}{按 $\sqrt{N_{\max}/N_c}$ 重复采样稀有害虫图片。}{bluelight}
\end{column}
\begin{column}{0.48\textwidth}
\carditem{两阶段损失}{RPN：前景/背景 BCE + 框回归 Smooth L1；RoI Heads：类别 CE + 框回归 Smooth L1。}{orangelight}
\vspace{2mm}
\carditem{推理后处理}{NMS IoU=0.5，每图最多 30 框；测试 AP 使用 score 阈值 0.05。}{greenlight}
\end{column}
\end{columns}
\end{frame}

% Detection history
\begin{frame}{检测模型训练过程}
\fig{0.95}{0.62}{assets/detection_history.png}
\vspace{-2mm}
\begin{card}{bluelight}
第 8 轮验证 mAP@0.5:0.95 最高（35.35\%），连续 6 轮未提升后早停；官方测试集仅评估一次。
\end{card}
\end{frame}

% Detection final metrics
\begin{frame}{检测模型最终测试结果}
\begin{columns}[T]
\begin{column}{0.5\textwidth}
\begin{table}\centering\small
\begin{tabular}{lll}
\toprule
\textbf{指标} & \textbf{结果} & \textbf{含义} \\
\midrule
mAP@0.5:0.95 & 34.3449\% & IoU 0.50$\sim$0.95 平均 AP \\
AP50 & 59.6817\% & IoU=0.5 时平均精度 \\
AP75 & 34.8622\% & IoU=0.75 时平均精度 \\
Precision@0.5 & 88.1562\% & 高分框中正确比例 \\
Recall@0.5 & 76.7102\% & 真实框被检出比例 \\
\bottomrule
\end{tabular}
\end{table}
\end{column}
\begin{column}{0.46\textwidth}
\fig{1.0}{0.46}{assets/detection_metrics.png}
\end{column}
\end{columns}
\end{frame}

% Key decisions
\begin{frame}{关键设计决策}
\begin{table}\centering\small
\begin{tabular}{p{0.34\textwidth}p{0.60\textwidth}}
\toprule
\textbf{决策} & \textbf{理由} \\
\midrule
分类 384，检测 640 & 分类保留纹理，检测保留定位精度 \\
检测骨干复用分类骨干 & DLCPD-25 已学到农业图像特征，迁移至害虫检测 \\
EMA 权重用于验证与保存 & 平滑训练后期波动，稳定 best 模型 \\
两个任务独立早停 & 避免一个任务过拟合影响另一个任务 \\
测试集只评估一次 & 严格隔离模型选择与最终评估 \\
\bottomrule
\end{tabular}
\end{table}
\end{frame}

% Inference flow
\begin{frame}{推理系统与 Web 应用}
\begin{center}
\small\bfseries
上传图片 \quad$\rightarrow$\quad EXIF 校正 / RGB \quad$\rightarrow$\quad 分类专家 Top-5 \quad$\rightarrow$\quad 检测专家害虫框 \quad$\rightarrow$\quad class\_id 映射 \quad$\rightarrow$\quad 标注图 + JSON
\end{center}
\vspace{3mm}
\begin{card}{greenlight}
FastAPI 提供 \texttt{POST /api/analyze}；返回分类 Top-5、检测框与 base64 标注图。训练进度看板每 3 秒刷新。
\end{card}
\end{frame}

% Engineering
\begin{frame}{工程实现与复现}
\begin{columns}[T]
\begin{column}{0.48\textwidth}
\carditem{分类训练}{python -m dlcpd25\_v2.classification.train}{greenlight}
\vspace{2mm}
\carditem{检测训练}{python -m dlcpd25\_v2.detection.train}{bluelight}
\vspace{2mm}
\carditem{最终评估}{python -m dlcpd25\_v2.classification.evaluate 或 detection.evaluate}{orangelight}
\vspace{2mm}
\carditem{启动 Web}{python -m dlcpd25\_v2.web --config configs/plan-a/app.yaml}{greenlight}
\end{column}
\begin{column}{0.48\textwidth}
\begin{tcolorbox}[colback=navy, colframe=navy, arc=2mm, top=2mm, bottom=2mm, left=3mm, right=3mm]
\color{white}\footnotesize
\texttt{project/src/dlcpd25\_v2/}\\[1mm]
\texttt{|-- classification/}\\[1mm]
\texttt{|-- detection/}\\[1mm]
\texttt{|-- serving/}\\[1mm]
\texttt{`-- web/}
\end{tcolorbox}
\end{column}
\end{columns}
\end{frame}

% Limitations
\begin{frame}{限制与未来工作}
\carditem{检测小目标仍弱}{AP small 仅 9.13\%，可引入 P2 增强、更高分辨率或 DETR/YOLO。}{orangelight}
\vspace{2mm}
\carditem{稀有类别样本少}{IP102 部分类仅数十框，可使用原始 75k 分类图做弱监督/伪标签。}{orangelight}
\vspace{2mm}
\carditem{分类 Macro-F1 有提升空间}{可尝试层级 softmax、更大 backbone、蒸馏或更强的尾类采样。}{greenlight}
\vspace{2mm}
\carditem{部署优化}{当前 PyTorch 权重较大，可导出 ONNX/TensorRT，检测可替换 YOLO 系列。}{bluelight}
\end{frame}

% Conclusion
{
\setbeamertemplate{footline}{}
\begin{frame}[plain]
\begin{tikzpicture}[remember picture,overlay]
  \fill[navy] (current page.south west) rectangle (current page.north east);
  \fill[green] ([yshift=-0.08\paperheight]current page.north west) rectangle ([yshift=-0.085\paperheight]current page.north east);
  \node[text=white,font=\LARGE\bfseries] at ([yshift=-0.26\paperheight]current page.north) {总结};
  \node[text=white,anchor=west,font=\normalsize] at ([xshift=0.12\paperwidth,yshift=-0.39\paperheight]current page.north west)
    {分类专家 ConvNeXt-Tiny@384：Top-1 89.22\%，Macro-F1 75.35\%};
  \node[text=white,anchor=west,font=\normalsize] at ([xshift=0.12\paperwidth,yshift=-0.47\paperheight]current page.north west)
    {检测专家 ConvNeXt-Tiny-FPN@640：mAP50:95 34.34\%，Precision 88.16\%};
  \node[text=white,anchor=west,font=\normalsize] at ([xshift=0.12\paperwidth,yshift=-0.55\paperheight]current page.north west)
    {双专家独立训练、独立评估、统一编排，推理链路完整可复现};
  \node[text=green,font=\Large\bfseries] at ([yshift=-0.78\paperheight]current page.north) {感谢观看 · 欢迎提问};
\end{tikzpicture}
\end{frame}
}

\end{document}
